\section{Summary and Research Gap}

The previous sections have reviewed the main approaches to code translation,
including traditional methods, deep learning-based techniques, and Transformer-
based models. These approaches have significantly improved the ability of models
to learn syntactic and semantic patterns from large-scale datasets.

However, several important challenges remain. One key limitation is that current
models often fail to preserve functional correctness, particularly in complex or real-
world scenarios. Additionally, the availability of high-quality parallel datasets is still
limited, which restricts the generalization capabilities of models.

Another major challenge lies in evaluation. Widely used metrics such as BLEU
and CodeBLEU focus primarily on surface-level similarity between generated
and reference code. While CodeBLEU incorporates syntactic and semantic features, it still does not directly evaluate functional correctness \cite{ren2020codebleu}.

Furthermore, studies have shown that these similarity-based metrics do not reli-
ably correlate with human judgment or actual correctness, highlighting the limita-
tions of current evaluation approaches \cite{evtikhiev2022bleu}.

These limitations reveal a clear research gap: there is a lack of reliable evaluation
methodologies that can accurately assess whether translated code is functionally
correct.

This thesis addresses this gap by focusing on execution-based evaluation. By
compiling and testing generated code against test cases, the proposed approach
aims to provide a more reliable assessment of semantic correctness compared to
traditional similarity-based metrics.

%SzB: Nem szoktak külön összefoglalót a related workhöz írni, meg nem így szokták összefoglalni a limitationöket
% Amit itt írok, a 2_6 és 2_7re egyaránt vonatkozik
% A chapter nagyjából ok, de ez a két utolsó rész be kell olvadjon a többibe
% Ahol bemutatsz valamit aminek vannak limitációi, ott lehet megemlíteni félmondatban
% Az is szokott lenni, hogy:
% X et al Y-t csinált, ami ilyen és olyan volt, ugyanakkor Z szempontból nem volt tökéletes
% A et al ezen javított azzal, hogy B-t használt. De persze ez se teljesen ok, mert C.
% És akkor ha itt hirtelen abbahagyod, azzal implicit sejteted, hogy a te munkád amúgy erre fog megoldást kínálni.  